\subsection{The full common-word audit for the Krammer test}

\paragraph{D23.P9 (all reduced pentagon assignments have finite-word realizations).}
Retain D23.P6--P8's notation, and write
$G_{43}=\langle X_4,X_3\rangle$ with the indicated reversed
marking. The simultaneous map
\[
 \widehat F'\longrightarrow
 P_{12}'\times G_{34}'\times G_{45}'\times G_{43}'
\tag{D23.P9.1}
\]
is onto. Every one of the $216$ assignments of D23.P7 can
therefore be realized by an ordinary free commutator word whose
five actual Krammer evaluations multiply to $I$ in the displayed
pentagon order. No unrelated-word substitution is involved.

This is a statement about ordinary finite word images. It does
not prove that one of these words is a gentle parameter. In
particular it does not construct an actual full-GT element with
nonzero $\Delta$.

\begin{proof}
Use D23.C9's faithful decompositions of each of the three mixed
pair groups into a solvable block of order $108$ and a unitary
block of order $24192$. On the reversed pair use the same basis
as pair $34$ and exchange its two generators. Their derived
groups are respectively $H_j'$ of order $27$ and simple
$S_j=\operatorname{SU}_3(3)$ of order $6048$.

The exact finite calculation
\path{agents/pentagon/krammer_allfive_commonword.py} forms the
three block-diagonal solvable generator pairs and computes the
normal closure of their joint commutator. It contains exactly
$19683=27^3$ matrices. The conjugacy orbit of that commutator
has $432$ elements and six of them generate the entire closure.
The image is contained in the product of the three order-$27$
derived groups, so this count proves equality with that full
product. This is a complete finite enumeration, not a random
test. Normal closure of the commutator computes the derived
subgroup of a two-generator group because the resulting
quotient makes the two generators commute.

The simple factors are independent as well. D23.P8's word
$w_1$ of length $28$ has identity pattern
\[
  (w_1(S_{34}),w_1(S_{45}),w_1(S_{43})):
       \quad(I,\ne I,\ne I).
\]
A second explicit word is
\[
 w_2=yxy^2xyx^{-1}y^{-3}x^{-1}y^{-1}
       xyx^2yxy^{-1}x^{-3}y^{-1}x^{-1}.
\tag{D23.P9.2}
\]
It has exponent sums zero and satisfies
$w_2(S_{45})=I$, $w_2(S_{43})\ne I$.
The program verifies these literal word values in the stored
three-dimensional matrices after restriction of scalars to
$\mathbb F_3$. Hence the kernels on $F'$ of the three simple
projections are pairwise distinct. A subdirect product of
nonabelian simple groups is a product of diagonal factors;
distinct kernels exclude every identification. The joint image
is therefore $S_{34}\times S_{45}\times S_{43}$.
The standard simplicity input and primary locators are those
recorded explicitly in D23.P8 and D23.F14.

By D23.P6, the abelianization of the solvable group $P_{12}'$
has order eight. The joint solvable mixed image is a three-group.
They have no nontrivial common quotient, because any nontrivial
finite solvable quotient has a nontrivial abelian quotient.
Goursat therefore separates $P_{12}'$ from the three solvable
mixed factors. A further application of Goursat separates this
solvable product from the perfect product of the three simple
factors. This proves (D23.P9.1).

Now take a listed solution $(F,R)$ of D23.P7 and put
\[
 H=Q_0^{-1}(C_0^{-1}RC_0)Q_0(C_0^{-1}R^{-1}C_0).
\]
Surjectivity provides a common commutator word $f$ with
\[
 f(X_1,X_2)=F,\quad f(X_3,X_4)=R,\quad
 f(X_4,X_5)=H,\quad f(X_4,X_3)=R^{-1}.
\]
The first-pair half twist $S=ABA$ exchanges $X_1,X_2$.
Thus, for every word, not merely for gentle words,
$f(X_2,X_1)=Sf(X_1,X_2)S^{-1}$. The selected six first-pair
matrices satisfy $SFS^{-1}=F^{-1}$ by D23.F13. The reversing
half twist $J=ABCABA$ gives
\[
 f(X_2,X_3)=JF^{-1}J^{-1},\qquad
 f(X_5,X_1)=JR^{-1}J^{-1}.
\]
The actual five values of this one word are consequently the
five factors in D23.P7.1. Their product is $I$ by the checked
assignment. The finite image of the ordinary free commutator
subgroup equals that of its profinite closure, so $f$ can be
chosen as an ordinary word.

The argument prescribes only finitely many evaluations and their
displayed relations. It does not establish the identities in
the full profinite free group required for a gentle pair, which
is the stated limitation.
\end{proof}
